\documentclass[11pt]{article}
\usepackage[margin=25mm]{geometry}
\usepackage{fontspec}
\setmainfont{Latin Modern Roman}
\usepackage{amsmath,amssymb,mathtools}
\usepackage{unicode-math}
\setmathfont{Latin Modern Math}
\usepackage{xurl,hyperref}
\hypersetup{colorlinks=true,urlcolor=blue}
\setcounter{secnumdepth}{0}
\setlength{\parindent}{0pt}
\setlength{\parskip}{5pt}
\setlength{\emergencystretch}{3em}
\newcommand{\sref}[2]{\hyperlink{source:#1:#2}{[S#2]}}
\newcommand{\cataloguescope}[1]{\paragraph{Recorded target.}#1}
\newcommand{\R}{\mathbb R}
\newcommand{\C}{\mathbb C}
\newcommand{\N}{\mathbb N}
\newcommand{\Z}{\mathbb Z}
\newcommand{\Q}{\mathbb Q}
\newcommand{\F}{\mathbb F}
\newcommand{\eps}{\varepsilon}
\begin{document}
% UnsolvedMath ID 30002801; source catalogue code OWR-13495-002
\setcounter{section}{30002800}
\section{Construction of K-Matrices for Coideal Subalgebras}\label{30002801}
{\small UnsolvedMath ID 30002801 · Mathematical Physics}
\subsection{Definitions and mathematical statement}

\paragraph{Shared notation.}$\mathbb N=\{1,2,\ldots\}$; the standard number fields have their usual meanings. All additional conventions are stated locally.


\paragraph{Definitions and conventions.}
An $R$-matrix encodes the exchange relation of an integrable model. A boundary $K$-matrix solves the corresponding reflection equation; one common additive-parameter convention is $R_{12}(u-v)K_1(u)R_{21}(u+v)K_2(v)=K_2(v)R_{12}(u+v)K_1(u)R_{21}(u-v)$. Subscripts mark tensor legs and $R_{21}=PR_{12}P$. A coideal subalgebra $B$ of a Hopf algebra $A$ has coproduct taking values in $B\otimes A$ or $A\otimes B$, according to convention.

\paragraph{Target retained from the canonical source catalogue.}
\paragraph{Catalogue wording retained for traceability.}
Develop a general method for constructing rational and trigonometric $K$-matrices satisfying the matrix reflection equations associated with twisted Yangians and quantum-symmetric-pair coideal subalgebras.

\paragraph{Source scope and qualification.}
The displayed equation only explains an ordinary reflection-equation convention. Twists, representations and multiplicative spectral parameters must be those of the report; both twisted Yangian and quantum-symmetric-pair cases remain in the target.
\subsection{Short English statement}
Find a general way to construct the rational and trigonometric boundary matrices for the indicated coideal algebras.
\subsection{Research attempt}
\paragraph{Model, date and investigated scope.}
OpenAI GPT-6 Astra (configured model \texttt{gpt-6-astra}), ultra reasoning; generated 11 October 2026. Verifies scalar rational and constant-square rational reflection-equation ansätze by direct tensor-leg algebra.

\paragraph{Sources and literature check.}
The complete supplied definition, including source qualifications and the dated literature assessment, was reviewed. The calculation below is an elementary restricted attempt motivated by [S2]; it does not claim novelty or a fresh literature-status determination.

\paragraph{Context and investigated scope.} The complete supplied definition, including source qualifications and the dated literature assessment, was reviewed. The calculation below is an elementary restricted attempt motivated by [S2]; it does not claim novelty or a fresh literature-status determination.

\paragraph{A rational scalar-boundary check.} On $V\otimes V$ let $P(x\otimes y)=y\otimes x$ and take $R(u)=uI+\eta P$, with fixed scalar $\eta$. Since $P^2=I$, all $R(u)$ commute and $R_{21}(u)=R(u)$. For any scalar function $k(u)$ the boundary matrix $K(u)=k(u)I_V$ satisfies the displayed ordinary reflection equation: both sides equal $k(u)k(v)R(u-v)R(u+v)$. This verifies a rational family, including normalization choices $k(0)=1$ if desired.

\paragraph{A less trivial constant ansatz.} Let $K$ be a constant matrix with $K^2=cI$. Put $A=K\otimes I$, $B=I\otimes K$, so $AB=BA$, $PA=BP$ and $PB=AP$. Expanding the two sides for $R(x)=xI+\eta P$, $x=u-v$, $y=u+v$, their difference reduces to $\eta x(A^2-B^2)P$, which vanishes under $K^2=cI$. The terms proportional to $xy$ and $\eta^2$ cancel using these identities, and the other mixed terms coincide. Singular $K$ is allowed algebraically, although a physical boundary may require invertibility.

\paragraph{Why this is not the requested method.} These are ordinary rational reflection-equation checks, with no coideal intertwining, twist or trigonometric parameter dependence. A solution of this matrix equation need not arise from the specified twisted Yangian or quantum symmetric pair. The first gap is relating representation-theoretic coideal data to a normalized spectral-parameter family in both source cases. The elementary ansätze do not settle that general construction problem.

\paragraph{Verification and adversarial audit.}
Checked permutation identities, expanded mixed terms, scalar normalization and singular constant K boundary; distinguished ordinary from twisted equations. Model self-audit only; no independent human proof checking. Prompts A and B were applied to this canonical record and its qualified source scope. Checked permutation identities, expanded mixed terms, scalar normalization and singular constant K boundary; distinguished ordinary from twisted equations. Retained unresolved status and zero original-target progress; Coideal intertwining and twisted/trigonometric spectral-parameter constructions are not supplied.

Coordinator checked the scalar rational reflection equation and the constant-boundary expansion yielding the squared-matrix criterion. This additional check was model-assisted, not independent human verification.

\paragraph{Tools, process and contributions.}
Python 3 was used for file preparation and structural checks; no numerical experiment or theorem prover was used. Contribution prompt A was followed by prompt B. GPT-6 Astra with ultra reasoning drafted the research and performed its stated self-audit; any additional coordinator contribution is identified in the verification and audit. A GPT-6-based Codex coordinator assembled the manuscript and checked repository integration; its exact variant and reasoning configuration were not exposed. No human mathematical contribution, independent human proof check or formal proof certification is supplied. Complete derivations and the consulted source references are included here; private reasoning transcripts are not published.

\paragraph{Final status and first remaining gap.}
\textbf{UNRESOLVED}. Coideal intertwining and twisted/trigonometric spectral-parameter constructions are not supplied.

\section{Completion Estimate}
\noindent\textbf{COMPLETION: 0\%}
The calculation is an auxiliary or restricted attempt and establishes no newly verified portion of the full source target. Coideal intertwining and twisted/trigonometric spectral-parameter constructions are not supplied.

\subsection{Sources}
\begin{itemize}
\item[\hypertarget{source:30002801:1}{S1}] ULAM.AI and the UnsolvedMath Contributors, UnsolvedMath, record 30002801 (OWR-13495-002). Catalogue curation under CC BY 4.0; definitions and English exposition added locally.
\url{https://huggingface.co/datasets/ulamai/UnsolvedMath}.
\item[\hypertarget{source:30002801:2}{S2}] Mini-Workshop: Coideal Subalgebras of Quantum Groups (2015), Oberwolfach Reports, DOI 10.4171/owr/2015/10, pp. 533--569. Primary report accessed 2026-10-11; relevant PDF pages 24.
\url{https://doi.org/10.4171/owr/2015/10}.
\item[\hypertarget{source:30002801:3}{S3}] Regelskis and Vlaar, reflection matrices (catalogue-listed context; not independently consulted in this preparation).
\url{https://arxiv.org/abs/1602.08471}.
\item[\hypertarget{source:30002801:4}{S4}] Recent K-matrix classification (catalogue-listed context; not independently consulted in this preparation).
\url{https://doi.org/10.1016/j.jpaa.2026.108210}.
\end{itemize}
\label{end:30002801}
\end{document}
